\section{Mercy, invitation and the table}
\subsection{The Entrance: the sinner who can still approach}
The Entrance's biblical basis, Psalm 130:3--4 (Vulgate 129), begins its question at the point where self-defense fails. If the Lord were to keep a strict account of iniquity, who could stand? The question does not single out an especially disreputable class. It makes the petitioner acknowledge a need shared with others. The next assertion, that forgiveness is with God, prevents recognition of guilt from becoming a reason to stop praying. The whole psalm moves from the depths through waiting to Israel's hope of redemption. Its petition is already an act of turning toward the one who can help.

Augustine notices the breadth of the question. It is not merely whether this particular sinner can stand: no conscience can safely rest its hope in its own innocence. Yet his explanation of forgiveness does not end in an indulgent refusal to judge. In Christ's sacrifice he finds the redemption which frees the guilty and warns them concerning the life still ahead. Mercy accompanies the forgiven person on the way. Later in the exposition, Christ's Resurrection becomes the ground for the body's hope: the Head has gone before those who still await their own resurrection.\footnote{Augustine, \work{Enarrationes in Psalmos} 129.2--5, 7--8; NPNF1 VIII, under Psalm CXXX.}

Bellarmine sharpens the pastoral balance. Penitence needs knowledge both of misery and of mercy. A person who does not perceive the first sees no need of a cure; a person ignorant of the second can despair of any cure. His image of the sinner unable to climb out of the well explains why acknowledgment of need is not itself the power of rescue. The desired help comes from God.\footnote{Robert Bellarmine, \work{Commentary on the Book of Psalms}, Ps 129, ad vv. 3--5, John O'Sullivan's abridged English translation (1866).} At the beginning of a Mass whose Gospel scrutinizes a guest, this psalm supplies an indispensable posture: judgment is heard by people asking for mercy, not by spectators congratulating themselves on another person's exclusion.

\subsection{The Collect: good action surrounded by grace}
\label{proper-collect}
The Collect, \textit{Tua nos, quaesumus, Domine, gratia}, concerns grace before and after the faithful's good action. Its approved English is not printed here. The prayer places neither the beginning nor the continuation of goodness outside God's help. Nevertheless, good action remains something to be done. A request for enabling grace is not an excuse to defer the task indefinitely.

Thomas's exposition of Psalm 22 distinguishes mercy that inspires the heart from mercy that helps bring the work to completion. Both are needed.\footnote{Thomas Aquinas, \work{Super Psalmo} 22, n. 2, on mercy preceding and following; Latin reportatio in the Corpus Thomisticum witness.} The similarity illuminates the Collect's subject without making his psalm exposition a commentary on this oration. The assembly's petition is apt both for someone needing courage to begin and for someone whose initial generosity has grown tired. It answers a real difficulty raised by the garment: the required change is serious, but the person called to change is not abandoned to solitary moral effort.

\subsection{Isaiah: a feast whose fullness defeats death}
Isaiah 25:6--10a promises more than food after scarcity. The Lord prepares a feast for all peoples on the mountain, removes the covering spread over the nations, defeats death and wipes away tears. The people's response recognizes the God for whom they have waited. The reading ends with the Lord's hand resting on the mountain. Its culmination is a secure divine presence, not the guests' ability to supply their own feast.

The surrounding text gives the promise its scale. Isaiah 24 describes judgment on an earth disordered by transgression and ends with the Lord reigning on Mount Zion and in Jerusalem. Isaiah 25:1--5 praises his wonderful deeds and his refuge for the poor and needy. After the appointed ending, the chapter turns to the humiliation of Moab. The banquet's universal invitation stands within this larger movement of deliverance and judgment. Neither the universal scope nor the severity of the context should disappear. The appointed portion centers the joy of God's saving presence; it does not say that pride has become harmless.

Thomas links this feast explicitly to Matthew's royal banquet. He reads the mountain through Christ's Passion, from which saving goods come, and also through judgment. He then distinguishes the banquet of the Church, the interior banquet of the soul and the solemn banquet of heaven. In the first he joins remembrance of the Passion to the sweetness of love and the nourishment it gives; the second concerns love, contemplation and growth; the third concerns perfected fullness.\footnote{Thomas Aquinas, \work{Expositio super Isaiam ad litteram}, cap. 25, paragraph 84985, exposition of \textit{Et faciet Dominus} and final \textit{Nota} on the threefold banquet.} These are connected forms of communion rather than interchangeable scenes. The distinction makes it possible to speak of the Eucharistic table as a present participation without pretending that death has already vanished from the congregation.

The imagery has also received a specifically initiatory reading. The \work{Mystagogical Catecheses} transmitted under Cyril of Jerusalem cites Isaiah's feast in its instruction on chrism. That exposition uses the Greek wine, gladness and anointing form of the passage, which differs from the Hebrew-derived wording appointed here. Its immediate subject is anointing, not an exclusively Eucharistic interpretation.\footnote{\work{Mystagogical Catechesis} III.7, also numbered \work{Catechetical Lecture} XXI.7, NPNF2 VII. The traditional attribution to Cyril is retained; this study does not resolve the authorship question.} Oil and gladness therefore have a real place in the passage's Christian reception, but this witness cannot be made to quote the present Lectionary's wording.

\subsection{The Responsorial Psalm: provision within danger}
Psalm 23 (Vulgate 22) answers Isaiah with a voice that knows the Lord as shepherd and host. Pasture, water and right paths lead to the valley and then to the table. The order matters: the shepherd's provision does not require the dangerous valley to be absent. Nor must the enemies disappear before the table can be prepared. Presence within danger makes confidence possible. The final dwelling gives the journey a destination.

Augustine lets the Church speak to Christ. Water becomes baptismal refreshment; right paths depend on the Shepherd's name rather than the sheep's merit; rod and staff train a person growing into spiritual maturity. The table is nourishment suited to that maturity, with joy that displaces former empty delights. Mercy extends beyond the present mortal life toward abiding with the Lord.\footnote{Augustine, \work{Enarrationes} 22.1--6; NPNF1 VIII, Psalm XXIII.} This account does not reduce pastoral guidance to pleasant feeling. The Lord's discipline belongs to his care, and the recipient's growth belongs to the meaning of nourishment.

Thomas offers several related table readings: divine instruction, the sacramental altar and the heavenly table. Scripture strengthens against temptation; the sacramental table sustains the faithful against their enemies; the final table belongs to the kingdom. He interprets the lack of nothing in relation to what leads to salvation and to the ultimate possession of God.\footnote{Thomas, \work{Super Psalmo} 22, nn. 1--2.} Together these explanations allow the psalm to answer both a sinner's need for formation and a sufferer's need for accompaniment. They do not turn the response into a guarantee of undisturbed circumstances. The appointed stanzas and response are distinguished in the inventory; the complete Douay psalm below is study text, not the approved refrain.

\subsection{Philippians: strength that welcomes another's help}
Paul's final words of thanks concern a real relationship. The appointed verses, Philippians 4:12--14, 19--20, speak of being brought low and of abundance, of hunger and satisfaction, of strength in Christ and of the Philippians' sharing in his tribulation. The intervening verses, 15--18, are not read at this Mass, but they explain the gift sent through Epaphroditus and Paul's interest in the fruit it bears for the donors. The language of acceptable sacrifice in that context helps explain why this is more than an exchange of commodities.

Verse 13 draws its meaning from the circumstances just named. Paul can endure the conditions of his mission in Christ's strength. He is not announcing mastery over every possible enterprise. Chrysostom emphasizes that contentment must be learned and practiced. Plenty requires virtue as well as poverty: it can produce indolence or pride, just as want can tempt a person toward evil. Christ's strengthening prevents Paul's account of endurance from becoming self-praise.\footnote{John Chrysostom, \work{Homilies on Philippians} XV, ad 4:10--14, NPNF1 XIII.}

The next sentence matters equally. Paul commends their fellowship in his affliction. Chrysostom hears him rejoin the donors just when his independence might have made them think their help useless. Giving benefits the giver through generosity, not because heaven has a cash price. Thomas makes the practical consequence explicit: the virtue of knowing how to bear want does not mean that assistance should be withheld.\footnote{Chrysostom, \work{Philippians} XV, ad 4:14--18; Thomas, \work{Super Philippenses} 4, lectio 2, paragraph 87844.} The second interpretation develops this double truth: God's help strengthens the recipient, and the neighbor's help remains a genuine good.

\subsection{The Acclamation: hearts able to recognize hope}
The acclamation adapts Ephesians 1:17--18 rather than reproducing two complete verses. Its subject is the illumination by which the heart perceives the hope of God's calling. The biblical prayer arises from thanksgiving for faith and love and continues toward God's power displayed in Christ's Resurrection and exaltation. Hope has both a source and a destination; it is more than optimism about the next turn of events.

Chrysostom notices the pairing of faith and love and the prayer for a deeper understanding of what God has done and will give. The same divine power that raised Christ draws estranged people to God. Thomas distinguishes the hope appropriate to the present journey from the rich and stable inheritance of the life to come.\footnote{Chrysostom, \work{Homilies on Ephesians} III, ad 1:15--22, NPNF1 XIII; Thomas, \work{Super Ephesios} 1, lectio 6, paragraph 87792.} Before hearing a parable that can be mistaken for either indiscriminate welcome or bare threat, the assembly asks for sight adequate to its calling. Understanding here serves response: the good promised must become desirable enough to reorder the heart.

\subsection{Matthew: rejection, gathering and inspection}
Jesus speaks in the Temple dispute that begins in Matthew 21:23. Chief priests and elders question his authority; subsequent parables expose a refusal to respond to God's work. Matthew 21:45 names chief priests and Pharisees as recognizing themselves in the preceding parables. The royal wedding then heightens the contrast between a lavish invitation and those who disregard or attack its messengers. Farm and business are Matthew's stated distractions. The wife's marriage and purchased oxen in Luke's banquet story are not additional details of this proclaimed text.

The king does not begin by inspecting garments. He prepares, sends and sends again. Refusal includes both indifference and murderous hostility; judgment follows the killing of his servants. Then the invitation goes to the roads, and bad and good alike are gathered. These are the movements common to both Gospel forms. The short form ends here. Its filled room is a remarkable reversal of the original refusal, but it is still a mixed gathering.

Only the long form continues with the king's inspection, the guest's silence, the binding and expulsion, and the final called/chosen saying. The narrative does not explain the guest's omission by telling us that garments were issued at the door. Its severity cannot be settled by importing that unestablished custom, or by imagining that the man is condemned for inability to afford fine clothing. Gregory and Chrysostom interpret the garment morally and ecclesially. Gregory identifies charity, Chrysostom the life and practice fitting to the gift received. Their reasoning receives its fullest treatment in the first interpretation.\footnote{Gregory the Great, \work{Homiliae in Evangelia} 38.7--14; Chrysostom, \work{Homilies on Matthew} LXIX.2, NPNF1 X.}

Their treatment of the city differs. Chrysostom connects the burning with the Roman destruction of Jerusalem; Gregory reads the city's punishment through the body and soul under judgment, with the king's armies as angelic ministers. These are different explanations of the image, even though both insist that grace calls for a changed life.\footnote{Chrysostom, \work{Matthew} LXIX.1; Gregory, \work{Homiliae} 38.5.} The narrative's immediate conflict concerns the leaders addressed by Jesus; the long form also judges someone inside the newly gathered room. Ancient collective polemic against Jewish people does not establish inherited ethnic guilt or authorize violence against a people. The Christian hearer cannot assign all the warning elsewhere and keep the invitation for himself.

\subsection{Offerings and the two Communion alternatives}
\label{proper-offerings}
The Prayer over the Offerings, \textit{Suscipe, Domine, fidelium preces}, concerns the faithful's offerings and prayers in relation to heavenly glory. The Latin witness controls this subject; an exact approved U.S. English exemplar remains unavailable, and no English prayer is supplied. The offering is directed beyond the immediate value of the thing brought. Paul's gift-context similarly distinguishes the material consumed from the good borne in those who give. That connection concerns the action's theological meaning, not an identity between Paul's collection and the Missal prayer.

The psalm Communion option adapts Psalm 34:11 (Vulgate 33:11). Seekers of the Lord are promised good, while the wealthy can lack. The whole psalm joins deliverance, praise, instruction in righteous conduct and acknowledgment of the just person's many afflictions. Augustine refuses to make wealth the measure of its truth: virtues and Christ the living bread are real riches that money cannot buy. Thomas also distinguishes spiritual good from perishable possessions and relates the promise to Paul's trained contentment.\footnote{Augustine, \work{Enarrationes} 33, second exposition, 13--14; Thomas, \work{Super Psalmo} 33, nn. 10--11.} The problem posed by an apparently prosperous wrongdoer is not a modern embarrassment added to the psalm. Augustine himself makes it central to his explanation.

The alternative Communion excerpts 1 John 3:2. The complete verse first establishes present divine childhood, then distinguishes what has not yet appeared and promises likeness through sight. Augustine's exposition makes that interval fruitful: desire stretches the capacity to receive God. The next verse's purification is an action of the believer with God's help. Likeness does not make creature and Creator equal.\footnote{Augustine, \work{Homilies on the First Epistle of John} IV.5--9, NPNF1 VII.} The psalm option concentrates on the good of seeking the Lord; the Johannine option names the vision toward which seeking moves. They are alternatives, not two consecutive antiphons.

\subsection{After Communion: nourishment for divine life}
\label{proper-after-communion}
The Prayer after Communion, \textit{Maiestatem tuam, Domine}, joins sacramental nourishment to participation in divine life. Its approved English is not reproduced. The prayer makes transformation, rather than mere completion of an action, the purpose of reception. Isaiah's final feast and John's future likeness give this purpose a horizon greater than present comfort.

Augustine's distinction between likeness and equality clarifies the kind of participation at stake. The Son's equality with the Father is not the creature's mode of sharing in God's goodness. The faithful are truly given life and formed for sight, while remaining creatures. His account of purification also makes participation morally active: the one being nourished is called to live differently. Thomas's distinction between the Church's banquet and heaven keeps the sacramental gift real without collapsing the journey into its completed end. The prayers have carried the assembly from dependence in action to a life received and still growing toward its fullness.
